\heading{理论物理基础（II）-课程说明}
\noteinfo{以下说明依据本项目现有真题、模拟题与参考答案整理，仅用于复习导航，不替代任课教师当学期发布的正式教学安排。}

\textbf{一、资料概况}

本项目当前收录本课程期中、期末真题各 1 套，期中模拟题 6 套，期末模拟题 6 套。整体看，这门课在资料中分成两大板块：\textbf{期中偏热学与统计物理基础，期末偏量子力学基础}。因此复习时不宜把全部内容混在一起刷题，最好按“期中范围 / 期末范围”分开推进。

\textbf{二、考试范围（按现有题库归纳）}

\textbf{期中常见范围：}
\begin{itemize}[leftmargin=2em,itemsep=0.2em,topsep=0.2em]
  \item 热力学第一、第二定律，热容、自由能、焓、Gibbs 自由能与 Maxwell 关系；
  \item 经典理想气体、非理想气体的状态方程与热力学函数；
  \item 正则系综、巨正则系综、配分函数与热力学量；
  \item 经典极限、Boltzmann 分布，以及玻色 / 费米分布的基本应用；
  \item 光子气体、量子统计、低温或高温极限下的近似计算。
\end{itemize}

\textbf{期末常见范围：}
\begin{itemize}[leftmargin=2em,itemsep=0.2em,topsep=0.2em]
  \item 波函数、表象、测量拟设、概率振幅与时间演化；
  \item 一维定态问题：无限深势阱、有限方势垒、$\delta$ 势、简谐振子等；
  \item 三维中心势、球谐函数、轨道角动量与升降算符；
  \item 简并、对称性、两能级系统、时间演化与序贯测量；
  \item 全同粒子、玻色子 / 费米子构造与简单多粒子体系。
\end{itemize}

\textbf{三、内容与命题特点}

这门课的特点是\textbf{跨度大但套路清楚}：期中更看重热统中的定义辨析、偏导变换与配分函数计算；期末更强调量子力学中“会不会写本征态、会不会展开、会不会算测量概率”。就现有题库而言，命题通常不故意追求偏题，而是要求你在标准模型中把\textbf{推导链条写完整}。

\textbf{四、难度评价}

整体难度可评为\textbf{中上}。单个知识点往往不算冷门，但一道题里常把多个环节串在一起，例如先求本征态再做展开，再讨论测量与时间演化；或者先由热力学关系推出函数形式，再进一步求化学势或稳定性判据。真正拉开差距的，通常不是记忆公式多少，而是\textbf{能否在有限时间内保持推导干净、变量不混乱、物理解释不跑题}。

\textbf{五、对课程的基本评价}

从现有试卷看，这门课非常适合通过题目建立知识框架。原因是它的典型题型高度代表课程核心：热统部分训练“从定义到函数关系”的组织能力，量子部分训练“从算符到本征展开再到测量”的表达能力。只要把每类标准模型吃透，复习收益会很高。

\textbf{六、如何更好利用模拟题}
\begin{itemize}[leftmargin=2em,itemsep=0.2em,topsep=0.2em]
  \item 先按板块刷题：先集中做完期中或期末，不建议冷热统与量子交替练习；
  \item 第一轮以“写全步骤”为主，第二轮再压时间；
  \item 对每一套题，至少总结一次“这题用了哪个标准模型、哪一步最容易错”；
  \item 若一道题涉及展开系数、概率、时间演化三步，务必整题重做，不要只背最后答案；
  \item 期中部分建议整理常见热力学恒等式与配分函数公式表，期末部分建议整理一页常见本征态与升降算符关系表。
\end{itemize}

